% Options for packages loaded elsewhere
\PassOptionsToPackage{unicode}{hyperref}
\PassOptionsToPackage{hyphens}{url}
\PassOptionsToPackage{dvipsnames,svgnames,x11names}{xcolor}
\documentclass[
  10pt,
  fontset=fandol]{ctexart}
\usepackage{xcolor}
\usepackage[margin=16mm]{geometry}
\usepackage{amsmath,amssymb}
\setcounter{secnumdepth}{-\maxdimen} % remove section numbering
\usepackage{iftex}
\ifPDFTeX
  \usepackage[T1]{fontenc}
  \usepackage[utf8]{inputenc}
  \usepackage{textcomp} % provide euro and other symbols
\else % if luatex or xetex
  \usepackage{unicode-math} % this also loads fontspec
  \defaultfontfeatures{Scale=MatchLowercase}
  \defaultfontfeatures[\rmfamily]{Ligatures=TeX,Scale=1}
\fi
\usepackage{lmodern}
\ifPDFTeX\else
  % xetex/luatex font selection
\fi
% Use upquote if available, for straight quotes in verbatim environments
\IfFileExists{upquote.sty}{\usepackage{upquote}}{}
\IfFileExists{microtype.sty}{% use microtype if available
  \usepackage[]{microtype}
  \UseMicrotypeSet[protrusion]{basicmath} % disable protrusion for tt fonts
}{}
\makeatletter
\@ifundefined{KOMAClassName}{% if non-KOMA class
  \IfFileExists{parskip.sty}{%
    \usepackage{parskip}
  }{% else
    \setlength{\parindent}{0pt}
    \setlength{\parskip}{6pt plus 2pt minus 1pt}}
}{% if KOMA class
  \KOMAoptions{parskip=half}}
\makeatother
\setlength{\emergencystretch}{3em} % prevent overfull lines
\providecommand{\tightlist}{%
  \setlength{\itemsep}{0pt}\setlength{\parskip}{0pt}}
\usepackage{amsmath,amssymb,mathtools,needspace}
\xeCJKDeclareCharClass{CJK}{"2032,"0400->"04FF,"2160->"216B,"2460->"2473,"25A0,"25C7,"25CB,"25CF}
\IfFontExistsTF{Arial Unicode MS}{\xeCJKsetup{AutoFallBack=true}\setCJKfallbackfamilyfont{\CJKrmdefault}{Arial Unicode MS}}{}
\setcounter{secnumdepth}{0}
\setlength{\emergencystretch}{2em}
\allowdisplaybreaks
\usepackage{bookmark}
\IfFileExists{xurl.sty}{\usepackage{xurl}}{} % add URL line breaks if available
\urlstyle{same}
\hypersetup{
  pdftitle={平方根法},
  colorlinks=true,
  linkcolor={Maroon},
  filecolor={Maroon},
  citecolor={Blue},
  urlcolor={Blue},
  pdfcreator={LaTeX via pandoc}}

\title{平方根法}
\author{}
\date{}

\begin{document}
\maketitle

\subsubsection{7.4.2 平方根法}\label{ux5e73ux65b9ux6839ux6cd5}

应用有限元法解结构力学问题，最后归结为求解线性方程组，这时系数矩阵大多具有对称正定性质。所谓平方根法，就是利用对称正定矩阵的三角分解而得到的求解对称正定方程组的一种有效方法，目前在计算机上广泛应用平方根法解此类方程组。

设 \(A\) 为对称阵，且 \(A\) 的所有顺序主子式均不为零。由定理 7.3 知，\(A\) 可唯一分解为式（7.4.1）的形式。

为了利用 \(A\) 的对称性，将 \(U\) 再分解，即

\[
U=\begin{pmatrix}
u_{11}&&&\\
&u_{22}&&\\
&&\ddots&\\
&&&u_{nn}
\end{pmatrix}
\begin{pmatrix}
1&\dfrac{u_{12}}{u_{11}}&\cdots&\dfrac{u_{1n}}{u_{11}}\\
&\ddots&\ddots&\vdots\\
&&\ddots&\dfrac{u_{n-1,n}}{u_{n-1,n-1}}\\
&&&1
\end{pmatrix}=DU_0,
\]

其中 \(D\) 为对角阵，\(U_0\) 为单位上三角阵。于是

\[A=LU=LDU_0.\tag{7.4.6}\]

\href{../../source-images/pdf-195.jpeg}{原页扫描}

又

\[A=A^{\mathrm T}=U_0^{\mathrm T}(DL^{\mathrm T}),\]

由分解的唯一性即得 \(U_0^{\mathrm T}=L\)，代入式（7.4.6）得到对称矩阵 \(A\) 的分解式 \(A=LDL^{\mathrm T}\)。

总结上述讨论，有以下定理。

\textbf{定理 7.7（对称阵的三角分解定理）} 设 \(A\) 为 \(n\) 阶对称阵，且 \(A\) 的所有顺序主子式均不为零，则 \(A\) 可唯一分解为

\[A=LDL^{\mathrm T},\]

其中 \(L\) 为单位下三角阵，\(D\) 为对角阵。

现设 \(A\) 为对称正定矩阵。首先说明 \(A\) 的分解式 \(A=LDL^{\mathrm T}\) 中 \(D\) 的对角元素 \(d_i\) 均为正数。事实上，由 \(A\) 的对称正定性，7.2 节的推论成立，即

\[d_1=D_1>0,\quad d_i=D_i/D_{i-1}>0\quad(i=2,3,\cdots,n).\]

于是

\[
D=\begin{pmatrix}d_1&&\\&\ddots&\\&&d_n\end{pmatrix}
=\begin{pmatrix}\sqrt{d_1}&&\\&\ddots&\\&&\sqrt{d_n}\end{pmatrix}
\begin{pmatrix}\sqrt{d_1}&&\\&\ddots&\\&&\sqrt{d_n}\end{pmatrix}
=D^{\frac12}D^{\frac12},
\]

由定理 7.7 得到

\[A=LDL^{\mathrm T}=LD^{\frac12}D^{\frac12}L^{\mathrm T}=(LD^{\frac12})(LD^{\frac12})^{\mathrm T}=L_1L_1^{\mathrm T},\]

其中 \(L_1=LD^{\frac12}\) 为下三角阵。

\textbf{定理 7.8（对称正定矩阵的三角分解或 Cholesky 分解）} 如果 \(A\) 为 \(n\) 阶对称正定矩阵，则存在一个实的非奇异下三角阵 \(L\) 使 \(A=LL^{\mathrm T}\)，当限定 \(L\) 的对角元素为正时，这种分解是唯一的。

下面用直接分解方法来确定计算 \(L\) 元素的递推公式。因为

\[
A=\begin{pmatrix}
l_{11}&&&\\
l_{21}&l_{22}&&\\
\vdots&\vdots&\ddots&\\
l_{n1}&l_{n2}&\cdots&l_{nn}
\end{pmatrix}
\begin{pmatrix}
l_{11}&l_{21}&\cdots&l_{n1}\\
&l_{22}&\cdots&l_{n2}\\
&&\ddots&\vdots\\
&&&l_{nn}
\end{pmatrix},
\]

其中 \(l_{ii}>0\)（\(i=1,2,\cdots,n\)）。由矩阵乘法及 \(l_{jk}=0\)（当 \(j<k\) 时），得

\[a_{ij}=\sum_{k=1}^{n}l_{ik}l_{jk}=\sum_{k=1}^{j-1}l_{ik}l_{jk}+l_{jj}l_{ij},\]

于是得到以下解对称正定方程组 \(Ax=b\) 的平方根法计算公式。

对于 \(j=1,2,\cdots,n\)，

\textbf{步 1}

\[l_{jj}=\left(a_{jj}-\sum_{k=1}^{j-1}l_{jk}^2\right)^{\frac12}.\tag{7.4.7}\]

\textbf{步 2}

\[l_{ij}=\left(a_{ij}-\sum_{k=1}^{j-1}l_{ik}l_{jk}\right)/l_{jj}\quad(i=j+1,\cdots,n),\]

\href{../../source-images/pdf-196.jpeg}{原页图像}

（续 7.4.2 平方根法。）

求解 \(Ax=b\)，即求解如下两个三角方程组：

（1）\(Ly=b\)，求 \(y\)；（2）\(L^{\mathrm T}x=y\)，求 \(x\)。

\[
\begin{aligned}
\text{步 3}\quad &y_i=\left(b_i-\sum_{k=1}^{i-1}l_{ik}y_k\right)/l_{ii}\quad(i=1,2,\cdots,n).\\
\text{步 4}\quad &x_i=\left(b_i-\sum_{k=i+1}^{n}l_{ki}x_k\right)/l_{ii}\quad(i=n,n-1,\cdots,1).
\end{aligned}
\tag{7.4.8}
\]

由式（7.4.7）知 \(a_{jj}=\sum_{k=1}^{j}l_{jk}^{2}\)（\(j=1,2,\cdots,n\)），所以

\[
l_{jk}^{2}\leqslant a_{jj}\leqslant\max_{1\leqslant j\leqslant n}\{a_{jj}\},
\]

\[
\max_{j,k}\{l_{jk}^{2}\}\leqslant\max_{1\leqslant j\leqslant n}\{a_{jj}\}.
\]

上面分析说明，分解过程中元素 \(l_{jk}\) 的数量级不会增长且对角元素 \(l_{jj}\) 恒为正数。于是不选主元素的平方根法是一个数值稳定的方法。

当求出 \(L\) 的第 \(j\) 列元素时，\(L^{\mathrm T}\) 的第 \(j\) 行元素亦就算出，所以平方根法约需 \(n^3/6\) 次乘除法运算，大约为一般直接 LU 分解法计算量的一半。

由于 \(A\) 为对称阵，因此在用计算机计算时只需存储 \(A\) 的下三角部分，共需要存储 \(n(n+1)/2\) 个元素，可用一维数组存放，即 \(\{a_{11},\allowbreak a_{21},\allowbreak a_{22},\allowbreak \cdots,\allowbreak a_{n1},\allowbreak a_{n2},\allowbreak \cdots,\allowbreak a_{nn}\}\)。

矩阵元素 \(a_{ij}\) 存放于一维数组的第 \(\left(i\cdot\frac{i-1}{2}+j\right)\) 个元素，\(L\) 的元素存放在 \(A\) 的相应位置。

由式（7.4.7）可看出，用平方根法解对称正定方程组时，计算 \(L\) 的元素 \(l_{ii}\) 需要用到开方运算，为了避免开方运算，下面用定理 7.7 的分解式 \(A=LDL^{\mathrm T}\)，即

\[
A=
\begin{pmatrix}
1&&&&\\
l_{21}&1&&&\\
l_{31}&l_{32}&1&&\\
\vdots&\vdots&\ddots&\ddots&\\
l_{n1}&l_{n2}&\cdots&l_{n,n-1}&1
\end{pmatrix}
\begin{pmatrix}
d_1&&&&\\
&d_2&&&\\
&&\ddots&&\\
&&&\ddots&\\
&&&&d_n
\end{pmatrix}
\begin{pmatrix}
1&l_{21}&l_{31}&\cdots&l_{n1}\\
&1&l_{32}&\cdots&l_{n2}\\
&&\ddots&\ddots&\vdots\\
&&&\ddots&l_{n-1,n-1}\\
&&&&1
\end{pmatrix},
\]

按行计算 \(L\) 的元素 \(l_{ij}\)（\(j=1,2,\cdots,i-1\)）。由矩阵乘法，并注意 \(l_{jj}=1\)，\(l_{jk}=0\)（\(j<k\)），得

\[
a_{ij}=\sum_{k=1}^{n}(LD)_{ik}(L^{\mathrm T})_{kj}
=\sum_{k=1}^{n}l_{ik}d_kl_{jk}
=\sum_{k=1}^{j-1}l_{ik}d_kl_{jk}+l_{ij}d_jl_{jj}.
\]

于是得到以下计算 \(L\) 的元素及 \(D\) 的对角元素公式：

对于 \(i=1,2,\cdots,n\)，

\[
\begin{aligned}
\text{步 1}\quad &l_{ij}=\left(a_{ij}-\sum_{k=1}^{j-1}l_{ik}d_kl_{jk}\right)/d_j\quad(j=1,2,\cdots,i-1).\\
\text{步 2}\quad &d_i=a_{ii}-\sum_{k=1}^{i-1}l_{ik}^{2}d_k.
\end{aligned}
\tag{7.4.9}
\]

\begin{quote}
校注（agent 补充，原书疑误）：式（7.4.8）步 4 原图确印 \(b_i\)，这里忠实保留。由本页所列 \(L^{\mathrm T}x=y\)，回代右端应使用 \(y_i\)；否则与其前方三角方程组不一致。\href{../assets/ch07b/196-step4.png}{原式局部图}。

校注（agent 补充，原书疑误）：\(LDL^{\mathrm T}\) 展开中最右矩阵次末行末项原图印作 \(l_{n-1,n-1}\)，这里保留原文；按左边 \(L\) 的转置，该位置应为 \(l_{n,n-1}\)。\href{../assets/ch07b/196-ldlt.png}{原矩阵局部图}。
\end{quote}

\href{../../source-images/pdf-197.jpeg}{原页图像}

（续 7.4.2 平方根法。）

为避免重复计算，引进 \(t_{ij}=l_{ij}d_j\)，由式（7.4.9）得到以下按行计算 \(L,T\) 元素的公式：

对于 \(i=1,2,\cdots,n\)，

\[
\begin{aligned}
\text{步 1}\quad &t_{ij}=a_{ij}-\sum_{k=1}^{j-1}t_{ik}l_{jk}\quad(j=1,2,\cdots,i-1).\\
\text{步 2}\quad &l_{ij}=t_{ij}/d_j\quad(j=1,2,\cdots,i-1).\\
\text{步 3}\quad &d_i=a_{ii}-\sum_{k=1}^{i-1}t_{ik}l_{ik}.
\end{aligned}
\tag{7.4.10}
\]

计算出 \(T=LD\) 的第 \(i\) 行元素 \(t_{ij}\)（\(j=1,2,\cdots,i-1\)）后，存放在 \(A\) 的第 \(i\) 行相应位置，然后再计算 \(L\) 的第 \(i\) 行元素，存放在 \(A\) 的第 \(i\) 行，\(D\) 的对角元素存放在 \(A\) 的相应位置。例如，

\[
A=\begin{pmatrix}
a_{11}&a_{21}&a_{31}&a_{41}\\
a_{21}&a_{22}&a_{32}&a_{42}\\
a_{31}&a_{32}&a_{33}&a_{43}\\
a_{41}&a_{42}&a_{43}&a_{44}
\end{pmatrix}
\longrightarrow
\begin{pmatrix}
d_1&&&\\
l_{21}&d_2&&\\
l_{31}&l_{32}&d_3&\\
t_{41}&t_{42}&t_{43}&t_{44}
\end{pmatrix}
\longrightarrow
\begin{pmatrix}
d_1&&&\\
l_{21}&d_2&&\\
l_{31}&l_{32}&d_3&\\
l_{41}&l_{42}&l_{43}&d_4
\end{pmatrix}.
\]

对称正定矩阵 \(A\) 按 \(LDL^{\mathrm T}\) 分解和按 \(LL^{\mathrm T}\) 分解计算量差不多，但 \(LDL^{\mathrm T}\) 分解不需要开方计算。求解 \(Ly=b,DL^{\mathrm T}x=y\) 的计算公式分别为步 4、步 5 所述公式。

\[
\begin{aligned}
\text{步 4}\quad &\begin{cases}
y_1=b_1,\\
y_i=b_i-\displaystyle\sum_{k=1}^{i-1}l_{ik}y_h\quad(i=2,\cdots,n).
\end{cases}\\
\text{步 5}\quad &\begin{cases}
x_n=y_n/d_n,\\
x_i=y_i/d_i-\displaystyle\sum_{k=i+1}^{n}l_{ki}x_k\quad(i=n-1,\cdots,2,1).
\end{cases}
\end{aligned}
\tag{7.4.11}
\]

计算公式（7.4.10）、（7.4.11）称为\textbf{改进的平方根法}。

\begin{center}\rule{0.5\linewidth}{0.5pt}\end{center}

\subsection{相关原页校注（agent 补充）}\label{ux76f8ux5173ux539fux9875ux6821ux6ce8agent-ux8865ux5145}

以下校注来自本节覆盖的原页，可能也涉及同页相邻小节；原文照录与校正说明分开保留。

\subsubsection{原书 PDF 页 197 的校注}\label{ux539fux4e66-pdf-ux9875-197-ux7684ux6821ux6ce8}

\href{../pages/pdf-197.md}{完整原页转录} · \href{../../source-images/pdf-197.jpeg}{原页图像}

校注记录：ch07b-E003。

\begin{quote}
校注（agent 补充，原书疑误）：式（7.4.11）步 4 的求和项原图印作 \(l_{ik}y_h\)，这里保留；求和指标为 \(k\) 且来源为 \(Ly=b\)，相应项应为 \(l_{ik}y_k\)。\href{../assets/ch07b/197-solves.png}{原式局部图}。
\end{quote}

\subsection{本节覆盖原页的校注记录}\label{ux672cux8282ux8986ux76d6ux539fux9875ux7684ux6821ux6ce8ux8bb0ux5f55}

下列记录按原页关联，含同页相邻小节的校注；计算时同时读取上文校注和完整原页。

\begin{itemize}
\tightlist
\item
  \texttt{ch07b-E001}：原书 PDF 页 196
\item
  \texttt{ch07b-E002}：原书 PDF 页 196
\item
  \texttt{ch07b-E003}：原书 PDF 页 197
\end{itemize}

\end{document}
